% %%%%%%%%%%%%%%%%%%%%%%%%%%%%%%%%%%%%%%%%%%%%%%%%%%%%%%%%%%%%%%%%%%%%%%%%%%%%%%%%
% 2345678901234567890123456789012345678901234567890123456789012345678901234567890
%         1         2         3         4         5         6         7         8

\documentclass[letterpaper, 10 pt, conference]{ieeeconf}  % Comment this line out if you need a4paper

%\documentclass[a4paper, 10pt, conference]{ieeeconf}      % Use this line for a4 paper

\IEEEoverridecommandlockouts                              % This command is only needed if
                                                          % you want to use the \thanks command

\overrideIEEEmargins                                      % Needed to meet printer requirements.

%In case you encounter the following error:
%Error 1010 The PDF file may be corrupt (unable to open PDF file) OR
%Error 1000 An error occurred while parsing a contents stream. Unable to analyze the PDF file.
%This is a known problem with pdfLaTeX conversion filter. The file cannot be opened with acrobat reader
%Please use one of the alternatives below to circumvent this error by uncommenting one or the other
%\pdfobjcompresslevel=0
%\pdfminorversion=4

% See the \addtolength command later in the file to balance the column lengths
% on the last page of the document

% The following packages can be found on http:\\www.ctan.org
%\usepackage{graphics} % for pdf, bitmapped graphics files
%\usepackage{epsfig} % for postscript graphics files
%\usepackage{mathptmx} % assumes new font selection scheme installed
%\usepackage{times} % assumes new font selection scheme installed
\usepackage{amsmath}  % assumes amsmath package installed
\usepackage{amssymb}  % assumes amsmath package installed
\usepackage{bm}       % bold math symbols

\title{\LARGE \bf
MAGIC: A Multi-Agent Dual-Arm System for \\ Automated Geometric Inspection and 3D Reconstruction
}

\author{Kartik Agrawal$^{*}$, Manyung Emma Hon$^{*}$, Kailash Jagadeesh$^{*}$, Megan Lee$^{*}$, and Shreya Ragi$^{*}$%
\thanks{*Authors contributed equally to this work.}%
\thanks{All authors are with the Robotics Institute, Carnegie Mellon University,
        Pittsburgh, PA 15213, USA. Project sponsored by Jiaoyang Li.
        {\tt\small \{kagrawal, mhon, kjagadee, meganl2, sragi\}@andrew.cmu.edu}}%
}


\begin{document}


\maketitle
\thispagestyle{empty}
\pagestyle{empty}


%%%%%%%%%%%%%%%%%%%%%%%%%%%%%%%%%%%%%%%%%%%%%%%%%%%%%%%%%%%%%%%%%%%%%%%%%%%%%%%%
\begin{abstract}

Automated inspection of complex industrial components, particularly those
produced via Wire Arc Additive Manufacturing (WAAM), presents significant
challenges due to self-occlusion, surface roughness, and geometric
irregularity. Traditional static sensor arrays or single-arm manipulator
systems often fail to achieve complete surface coverage of such parts. This
paper presents MAGIC (Multi-Agent Geometric Inspection and Classification), a
dual-arm robotic system designed for object-centered active inspection. Unlike
camera-in-hand approaches, MAGIC utilizes two coordinated 7-DOF manipulators
to reorient the workpiece in front of a high-precision 3D sensor array. We
detail the system's modular architecture, which integrates learning-based
handle detection, bimanual grasp planning, and Jacobian-based impedance
control to minimize occlusion. The motion planning pipeline is developed and
validated in a high-fidelity MuJoCo simulation of a Kinova Gen3 dual-arm
workcell before deployment on physical hardware. Experimental results
demonstrate that the system achieves a reconstruction completeness of over
90\% with a dimensional accuracy of 1.3\,cm RMSE, significantly
outperforming static single-view baselines.

\end{abstract}


%%%%%%%%%%%%%%%%%%%%%%%%%%%%%%%%%%%%%%%%%%%%%%%%%%%%%%%%%%%%%%%%%%%%%%%%%%%%%%%%
\section{INTRODUCTION}

The advent of Industry 4.0 has necessitated the integration of intelligent
inspection systems within manufacturing cells. Additive Manufacturing (AM)
techniques, such as Wire Arc Additive Manufacturing (WAAM), allow for the
creation of large-scale metal components with complex internal geometries
\cite{waam_overview}. However, the layer-by-layer deposition process in WAAM
often introduces geometric inconsistencies, surface voids, and significant
roughness that require rigorous quality assurance (QA) \cite{son2024coordinate}.

Current automated QA solutions largely rely on static 3D scanners or single
robotic arms equipped with wrist-mounted cameras. While effective for simple
convex geometries, these methods suffer from ``blind spots'' when inspecting
objects with deep concavities or undercuts. To scan the underside of a heavy
WAAM part, human intervention is typically required to manually flip the
object, introducing safety risks and breaking the digital thread of the
manufacturing process.

To address these limitations, we propose a ``system-centric'' or
``object-centered'' inspection paradigm. Instead of moving the sensor around
a stationary object, we propose manipulating the object itself to present
optimal views to a fixed sensor. This inverts the conventional next-best-view
paradigm and enables complete coverage without exposing human operators to
heavy, hot, or sharp WAAM components.

This paper makes the following contributions:

\begin{itemize}
    \item An object-centered dual-arm inspection paradigm that eliminates the
    need for camera-in-hand scanning for heavy industrial components.

    \item A Jacobian-transpose impedance control formulation for closed-chain
    rigid-body reorientation using synchronized dual manipulators, validated
    in a high-fidelity MuJoCo simulation.

    \item A learning-based multi-handle geometric pose estimation strategy
    for robust orientation recovery without fiducial markers.

    \item Experimental validation demonstrating over 90\% reconstruction
    completeness on occluded WAAM-proxy geometries.
\end{itemize}


%%%%%%%%%%%%%%%%%%%%%%%%%%%%%%%%%%%%%%%%%%%%%%%%%%%%%%%%%%%%%%%%%%%%%%%%%%%%%%%%
\section{RELATED WORK}

\subsection{Automated 3D Inspection}
Robotic inspection has been widely studied for automotive and aerospace
applications. Conventional approaches employ a ``Next-Best-View'' (NBV)
algorithm where a single arm moves a camera to minimize entropy in the
volumetric map \cite{bircher2016receding}. However, determining NBV in highly
constrained environments with a single manipulator often leads to unreachable
viewpoints or severe occlusion.

Son et al.\ \cite{son2024coordinate} highlighted the difficulties in
coordinate system setting for post-machining WAAM parts, emphasizing the need
for high-fidelity global surface models that static systems struggle to
provide due to specular reflectance and complex overhangs. MAGIC directly
addresses this by reorienting the workpiece itself.

\subsection{Bimanual Manipulation and Impedance Control}
Bimanual manipulation expands the robot's effective workspace and payload
capacity \cite{smith2012dual}. Smith et al.\ demonstrated collaborative
carrying tasks, but applied dual-arm coordination to logistics rather than
inspection. In the inspection domain, dual-arm systems have primarily been
used in ``master-slave'' configurations where one arm holds a light source and
the other a camera.

Our work differs fundamentally: the arms manipulate the \emph{workpiece},
requiring tight synchronization to prevent internal stresses on the rigid
object during rotation. Impedance control is the natural choice for
compliant interaction with the environment \cite{hogan1985impedance}. The
controller regulates the dynamic relationship between end-effector force and
position, allowing the arms to absorb compliance errors without damaging the
object or losing the grasp. We adopt a Cartesian-space impedance formulation
where the control torques are computed via the Jacobian transpose, avoiding
the need for explicit inverse dynamics.

\subsection{Point Cloud Registration}
Fusing multiple views into a single model is a classic computer vision
problem. Iterative Closest Point (ICP) and its variants are standard; however,
they are prone to local minima without good initialization
\cite{besl1992method}. By leveraging the high-precision forward kinematics of
the robot arms, MAGIC provides an accurate prior for each sensor-to-object
transform, significantly robustifying the multi-view registration process
compared to purely vision-based global registration methods.


%%%%%%%%%%%%%%%%%%%%%%%%%%%%%%%%%%%%%%%%%%%%%%%%%%%%%%%%%%%%%%%%%%%%%%%%%%%%%%%%
\section{SYSTEM ARCHITECTURE}

\subsection{Hardware Configuration}
The MAGIC workcell is constructed around two Kinova Gen3 7-DOF robotic
manipulators mounted on a shared, rigid Vention aluminum extrusion table. The
7 degrees of freedom are critical for redundant kinematic resolution, allowing
the arms to reorient objects without reaching joint limits or singularities.
Each arm is equipped with a Robotiq 85 two-finger parallel-jaw gripper,
providing a maximum grip force of 235\,N and a stroke of 85\,mm. The system
has a total of 16 actuated degrees of freedom: 7 per arm plus 1 finger joint
per gripper (the remaining gripper joints are coupled via tendon constraints).

The perception system comprises an Intel RealSense D435i RGB-D camera mounted
on a static overhead gantry. The D435i provides depth images at up to
90\,fps with a range accuracy of less than 2\% at 2\,m. The camera is
calibrated extrinsically to the robot base frame using a standard ChArUco
board routine, yielding a reprojection error of less than 1.5 pixels.

\subsection{Software Framework}

\subsubsection{Simulation Environment}
Motion planning and control algorithms are developed and validated in a
high-fidelity simulation environment built with MuJoCo \cite{todorov2012mujoco}.
The simulation model (Fig.~\ref{fig:architecture}) replicates the physical
Vention workcell geometry, including the two Kinova Gen3 arms, their Robotiq
85 grippers, and the table structure, using manufacturer-provided STL
meshes. The MuJoCo physics engine provides accurate rigid-body dynamics with
a timestep of 1\,ms, enabling realistic evaluation of impedance control gains
and closed-chain trajectory quality prior to hardware deployment.

\subsubsection{Deployment Architecture}
The deployed system targets ROS\,2 Humble to leverage its real-time DDS
communication layer. The architecture is divided into three primary
subsystems:

\begin{itemize}
    \item \textbf{Perception Node:} Handles data acquisition, object pose
    detection via the learning-based pipeline, and point cloud filtering
    (statistical outlier removal and voxel grid downsampling).

    \item \textbf{Control Node:} Implements the Cartesian-space impedance
    controller described in Section~\ref{sec:motion}, interfacing with the
    Kinova ROS\,2 driver to command joint torques at 1\,kHz.

    \item \textbf{Executive Node:} A finite state machine (FSM) that
    orchestrates the workflow, handling transitions between \textsc{Scan},
    \textsc{Grasp}, \textsc{Reorient}, and \textsc{Fuse} states.
\end{itemize}

\begin{figure}[thpb]
   \centering
   % \includegraphics[scale=1.0]{system_arch.png}
   \framebox{\parbox{3in}{}}
   \caption{The MuJoCo simulation model of the MAGIC workcell, showing the
   two Kinova Gen3 arms on the Vention table alongside the software
   architecture depicting asynchronous communication between the perception
   pipeline and the impedance control loop.}
   \label{fig:architecture}
\end{figure}


%%%%%%%%%%%%%%%%%%%%%%%%%%%%%%%%%%%%%%%%%%%%%%%%%%%%%%%%%%%%%%%%%%%%%%%%%%%%%%%%
\section{METHODOLOGY}
\label{sec:methodology}

The MAGIC inspection pipeline follows a sequential execute-and-verify
approach, defined by the following stages: (1)~Localization, (2)~Grasp
Synthesis, (3)~Coordinated Manipulation, and (4)~Reconstruction.

\subsection{Object Localization via Learning-Based Pose Estimation}

MAGIC employs a learning-based perception pipeline for real-time object pose
estimation rather than fiducial markers, which are impractical on rough WAAM
surfaces. A custom-trained YOLOv11 detector identifies four semantically
distinct handle classes on the workpiece. The handles are intentionally
color-coded and spatially arranged to encode geometric constraints that enable
stable, unambiguous orientation recovery.

Let the detected handle centers in image coordinates be $\{h_i\}_{i=1}^{4}$,
where each $h_i = (u_i, v_i)$ represents the pixel location of a detected
bounding box centroid. Using the camera intrinsic matrix $K$ and the aligned
depth map $D(u,v)$, each handle center is lifted into 3D camera coordinates:

\begin{equation}
    \bm{P}_i^{cam} = D(u_i, v_i) \, K^{-1}
    \begin{bmatrix} u_i \\ v_i \\ 1 \end{bmatrix}
\end{equation}

The object center is computed as the centroid of the four 3D handle positions:

\begin{equation}
    \bm{P}_{obj}^{cam} = \frac{1}{4} \sum_{i=1}^{4} \bm{P}_i^{cam}
\end{equation}

To compute the yaw orientation, we exploit the fixed geometric relationship
between a canonical pair of opposing handles (along the primary axis of the
part). Let $\bm{P}_a$ and $\bm{P}_b$ denote this pair. The yaw angle
$\theta$ about the vertical axis is:

\begin{equation}
    \theta = \operatorname{atan2}\!\left(P_b^y - P_a^y,\; P_b^x - P_a^x\right)
\end{equation}

The final object pose in the robot base frame is obtained via the known
extrinsic calibration transform $T_{base}^{cam}$:

\begin{equation}
    \bm{P}_{obj}^{base} = T_{base}^{cam}
    \begin{bmatrix} \bm{P}_{obj}^{cam} \\ 1 \end{bmatrix}
\end{equation}

This multi-handle geometric encoding eliminates the orientation ambiguity
commonly observed in single-marker or symmetric-object localization. Because
the pose is derived from the relative spatial configuration of multiple
detected features, the estimate is resilient to transient misdetections and
reduces frame-to-frame jitter. The resulting $(x, y, z, \theta)$ estimate is
published as a \texttt{PoseStamped} message, enabling closed-loop integration
with the manipulation subsystem.

\subsection{Grasp Synthesis and Stability}

For bimanual manipulation, selecting stable, collision-free grasp pairs is
critical. We simplify the problem by assuming the workpiece has two parallel
gripping surfaces separated by a known width $w$ (common in WAAM base plates).
The grasp planner generates candidate end-effector pose pairs $(g_L, g_R)$
for the left and right arms respectively.

A grasp pair is considered valid if all three conditions hold:

\begin{enumerate}
    \item Both $g_L$ and $g_R$ are reachable via numerical inverse
    kinematics (IK) within the joint limits of each 7-DOF arm.
    \item The Euclidean distance $\|g_L - g_R\|$ matches the object
    width $w$ within a tolerance $\epsilon = 5$\,mm.
    \item The grasp approach vectors are antipodal (anti-parallel) to
    ensure force closure on the workpiece.
\end{enumerate}

Valid pairs are ranked by a manipulability measure
$w(\bm{q}) = \sqrt{\det(J(\bm{q})\,J(\bm{q})^T)}$, and the highest-ranked
pair is selected. This maximizes distance from kinematic singularities
throughout the subsequent reorientation trajectory.

\subsection{Coordinated Motion via Impedance Control}
\label{sec:motion}

Reorienting a rigid object with two arms requires a closed-chain kinematic
constraint. If the arms move asynchronously, they exert damaging internal
forces on the workpiece or lose grip. We model the two arms and the
rigidly-grasped object as a single closed kinematic chain with
$n = 14$ degrees of freedom ($q \in \mathbb{R}^{14}$).

The constraint requires that for any time $t$ along the reorientation
trajectory, the relative transform between the two end-effectors remains
constant:

\begin{equation}
    ^{L}T_{R}(t) = {^{L}T_{R}}(0) \qquad \forall\, t \in [0, T]
    \label{eq:constraint}
\end{equation}

\subsubsection{Impedance Control Formulation}

We enforce constraint~\eqref{eq:constraint} in a decentralized fashion using
Cartesian-space impedance control on each arm independently, relying on the
high stiffness gains to maintain the internal constraint error below a safe
threshold. For each arm, the Cartesian position and orientation error
$\bm{e} \in \mathbb{R}^6$ is defined as the difference between the desired
and current end-effector pose (3D position + axis-angle orientation). The
corresponding joint torque command is:

\begin{equation}
    \bm{\tau} = J(\bm{q})^T \!\left(K_d\,\bm{e} + D_d\,\dot{\bm{e}}\right)
                + \bm{\tau}_{grav}(\bm{q})
    \label{eq:impedance}
\end{equation}

where $J(\bm{q}) \in \mathbb{R}^{6 \times 7}$ is the geometric Jacobian of
the end-effector, $K_d$ and $D_d$ are the stiffness and damping matrices, and
$\bm{\tau}_{grav}$ is the gravity-compensation torque (obtained from MuJoCo's
\texttt{qfrc\_bias}). The controller runs at 1\,kHz in simulation. The
stiffness and damping gains are set to:

\begin{equation}
    K_d = \operatorname{diag}(60,\,60,\,60,\,13,\,13,\,13)\;\text{N/m, Nm/rad}
\end{equation}
\begin{equation}
    D_d = \operatorname{diag}(4,\,4,\,4,\,2,\,2,\,2)\;\text{Ns/m, Nms/rad}
\end{equation}

Translational stiffness (60\,N/m) is set significantly higher than rotational
stiffness (13\,Nm/rad) to prioritize positional accuracy during reorientation
while allowing slight orientational compliance to absorb manufacturing
tolerances in the grasp tabs. To protect the joints and the object, joint
torque commands are clipped to $\pm 10$\,Nm. Gravity compensation ensures that
this torque budget is not consumed by the arm's static weight.

\subsubsection{Synchronized Trajectory Generation}

Desired end-effector trajectories for each arm are generated offline as a
pair of synchronized Cartesian waypoints that jointly rotate the object to the
target pose. Given the object center $\bm{c}$ and target rotation $R_{goal}$,
the desired pose of each arm's end-effector at each waypoint is computed
geometrically from the fixed grasp offsets. Waypoints are interpolated with a
minimum-jerk profile to keep internal forces within safe limits during
rotation.

\begin{figure}[thpb]
   \centering
   \framebox{\parbox{3in}{}}
   \caption{Simulated coordinated motion trajectory. The red and blue traces
   show the left and right end-effector paths respectively while maintaining
   the closed-chain constraint. The object undergoes a 90$^\circ$ rotation
   about its long axis.}
   \label{fig:motion_plan}
\end{figure}

\subsection{3D Reconstruction Pipeline}

The reconstruction module integrates depth frames captured at each stable
object pose into a global point cloud.

\subsubsection{Pre-processing}
Raw depth maps from the RealSense D435i are converted to point clouds in the
camera frame. We apply a Statistical Outlier Removal (SOR) filter
($k = 50$, $\sigma = 1.0$) to remove sensor noise and ``flying pixels''
common at the edges of metallic or specular objects.

\subsubsection{Registration}
As the robot reorients the object to pose $i$, we capture point cloud
$C_i$. The forward kinematics of the robot arm provide a high-accuracy
estimate of the sensor-to-object transform $T_{init}^{(i)}$. This serves
as the initial alignment guess for ICP, which refines the alignment by
minimizing the point-to-plane error:

\begin{equation}
    E = \sum_{\bm{p} \in C_i}
        \left\lVert \bigl(R\,\bm{p} + \bm{t} - \bm{q}\bigr)
                    \cdot \hat{\bm{n}}_q \right\rVert^2
\end{equation}

where $\bm{q}$ is the nearest neighbor of $R\,\bm{p} + \bm{t}$ in the
accumulated global model and $\hat{\bm{n}}_q$ is its estimated surface
normal.

\subsubsection{Fusion}
The ICP-aligned clouds are merged into a global model. A voxel grid with
leaf size 2\,mm is applied to normalize point density, removing duplicate
points in overlapping regions and producing the final mesh $M_{final}$.


%%%%%%%%%%%%%%%%%%%%%%%%%%%%%%%%%%%%%%%%%%%%%%%%%%%%%%%%%%%%%%%%%%%%%%%%%%%%%%%%
\section{EXPERIMENTAL EVALUATION}

\subsection{Experimental Setup}

\textbf{Simulation validation.}\; Motion planning and impedance control are
evaluated entirely in MuJoCo using the high-fidelity dual-arm simulation
described in Section~III.B. The simulation accurately models joint-level
dynamics, gravity compensation, and gripper-object contact, providing a
principled environment for gain tuning and trajectory quality assessment
before hardware trials.

\textbf{Reconstruction evaluation.}\; Experiments on 3D reconstruction are
conducted using 3D-printed mockups of WAAM impellers and brackets. These
proxies replicate the geometry of real parts but reduce the risk of damage
during development. Ground truth is defined by the CAD models used for
printing, which provide a reference surface with negligible manufacturing
error relative to the system's measurement precision.

\subsection{Metrics}

We evaluate the system on three key performance indicators (KPIs):

\begin{itemize}
    \item \textbf{Reconstruction Completeness ($C_p$):} The ratio of the
    surface area covered by the reconstructed scan to the total CAD surface
    area, expressed as a percentage.

    \item \textbf{Geometric Accuracy (RMSE):} The root mean square of the
    nearest-neighbor Euclidean distance from each reconstructed point to the
    CAD reference surface.

    \item \textbf{Cycle Time:} Total elapsed time from system initialization
    to final model generation.
\end{itemize}

\subsection{Results}

\subsubsection{Reconstruction Fidelity}

Table~\ref{table_results} compares MAGIC against a single-view static scan
baseline across three object classes. The single-view approach consistently
failed to capture the underside and internal faces, capping completeness at
45--58\%. MAGIC, through active reorientation to four canonical poses
(0$^\circ$, 90$^\circ$, 180$^\circ$, 270$^\circ$ about the vertical axis,
plus one lateral flip), achieved over 90\% completeness in all tested
geometries.

\begin{table}[h]
\caption{Reconstruction Performance Comparison}
\label{table_results}
\begin{center}
\begin{tabular}{|l||c|c|c|}
\hline
Object Type & Method & Completeness (\%) & RMSE (cm) \\
\hline
Simple Cube     & Static & 58.2 & 0.9 \\
Simple Cube     & MAGIC  & \textbf{98.5} & 1.1 \\
\hline
WAAM Impeller   & Static & 45.4 & 1.5 \\
WAAM Impeller   & MAGIC  & \textbf{90.1} & 1.3 \\
\hline
Overhang Part   & Static & 52.0 & 1.4 \\
Overhang Part   & MAGIC  & \textbf{89.4} & 1.6 \\
\hline
\end{tabular}
\end{center}
\end{table}

The slight increase in RMSE for MAGIC relative to single-view static scans is
attributed to error propagation during multi-view registration: each ICP step
introduces a small residual alignment error that accumulates across views.
Nevertheless, a worst-case RMSE of 1.6\,cm is acceptable for dimensional
verification of WAAM pre-forms prior to finish machining, where subsequent CNC
operations remove 3--5\,mm of material.

\subsubsection{Impedance Control Performance}

Closed-chain constraint satisfaction was evaluated in simulation over 50
reorientation trials. The mean constraint violation $\|{^{L}T_R(t)} -
{^{L}T_R(0)}\|_F$ remained below 3\,mm throughout all trials, confirming that
the impedance gains are sufficient to maintain the object without damaging
internal forces. The average time to complete a 90$^\circ$ reorientation
(including settling) was $4.4 \pm 3.4$\,s. The high standard deviation
reflects the sensitivity of ICP-driven settling detection to point cloud
density.

\subsubsection{System Latency}

The dominant computational bottleneck is the ICP registration step, which
scales linearly with point cloud density and contributes approximately
6--12\,s per captured view (at a voxel resolution of 2\,mm). Pose estimation
via YOLOv11 and handle lifting runs at real-time rates ($>$20\,Hz). Grasp
generation is near-instantaneous ($<$0.5\,s).

\begin{figure}[thpb]
   \centering
   \framebox{\parbox{3in}{}}
   \caption{Qualitative reconstruction comparison. (Left) Single static view
   showing large unscanned holes on the concave side of the WAAM impeller
   proxy. (Right) Fused multi-view model from MAGIC showing near-complete
   geometry coverage.}
   \label{fig:qualitative}
\end{figure}

\subsection{Grasp Reliability}

To test robustness, we executed 20 automated grasp-lift-replace cycles on the
physical workcell. The system achieved a 90\% (18/20) success rate. The two
failures were due to detection noise causing a 6--8\,mm misalignment of the
gripper fingers relative to the grasp tabs, leading to a slip during the
lift phase. Tightening the YOLOv11 confidence threshold from 0.5 to 0.7
eliminated these failures in subsequent offline tests but reduced the
detection rate for partially occluded handles.


%%%%%%%%%%%%%%%%%%%%%%%%%%%%%%%%%%%%%%%%%%%%%%%%%%%%%%%%%%%%%%%%%%%%%%%%%%%%%%%%
\section{DISCUSSION}

\subsection{Implications for Industry}

The ability to achieve near-complete reconstruction without human intervention
enables ``Lights-Out'' manufacturing inspection. The 3D models generated by
MAGIC can be directly fed into CAM software to generate hybrid manufacturing
toolpaths, where a CNC mill machines down rough WAAM surfaces to net shape.
The learning-based pose estimation approach avoids the need to attach or align
fiducial targets to hot or rough WAAM build plates, which is impractical in
production environments.

\subsection{Limitations}

\begin{itemize}
    \item \textbf{Sensor limitations:} The fixed overhead camera cannot
    resolve features inside deep, narrow channels even with object rotation.
    Augmenting the fixed sensor with a wrist-mounted camera on one arm would
    extend coverage to internal channels.

    \item \textbf{Processing speed:} The robot must pause motion at each
    view to capture stable frames and run ICP, limiting throughput.
    Continuous-motion scanning with online point cloud integration (e.g., via
    TSDF fusion) would remove this bottleneck.

    \item \textbf{Part rigidity:} The closed-chain rigid-body assumption holds
    for solid metal parts but fails for flexible or hollow-walled components,
    where dual-arm tension could deform the object.

    \item \textbf{Grasp generalization:} The current grasp planner assumes
    parallel gripping surfaces. Arbitrary WAAM geometries without dedicated
    grasp tabs require a learning-based grasp synthesis approach.
\end{itemize}


%%%%%%%%%%%%%%%%%%%%%%%%%%%%%%%%%%%%%%%%%%%%%%%%%%%%%%%%%%%%%%%%%%%%%%%%%%%%%%%%
\section{CONCLUSIONS AND FUTURE WORK}

We have presented MAGIC, a comprehensive dual-arm system for the automated
geometric inspection of complex industrial parts. By integrating Jacobian-based
impedance control for synchronized bimanual reorientation with a multi-handle
learning-based pose estimator and a multi-view 3D reconstruction pipeline, we
demonstrated a robust inspection workflow that resolves the occlusion problems
inherent in static single-view inspection cells. Motion planning and control
were developed and validated in a high-fidelity MuJoCo simulation before
physical hardware deployment.

Future work will focus on three areas:

\begin{enumerate}
    \item \textbf{Continuous-motion scanning:} Integrating TSDF-based volumetric
    fusion to allow point cloud capture during arm motion, removing the
    stop-and-scan latency.

    \item \textbf{Learning-based grasp synthesis:} Replacing the parallel-surface
    grasp planner with a 6-DOF grasp network to handle arbitrary WAAM
    geometries without dedicated grasp tabs.

    \item \textbf{Defect-aware toolpath generation:} Developing automatic
    M-code generation based on detected geometric deviations in the MAGIC
    output model, closing the loop with downstream CNC machining.
\end{enumerate}


%%%%%%%%%%%%%%%%%%%%%%%%%%%%%%%%%%%%%%%%%%%%%%%%%%%%%%%%%%%%%%%%%%%%%%%%%%%%%%%%
\section*{ACKNOWLEDGMENT}

The authors acknowledge the support of the Carnegie Mellon University Robotics
Institute and the guidance of our sponsor, Jiaoyang Li.

\bibliographystyle{IEEEtran}
\bibliography{references}


\end{document}
